\chapter{Acute Stress Disorder}
\index{Acute Stress Disorder}
\label{sec:Acute Stress Disorder}

\section{Overview}
Acute stress disorder has a variable prevalence (5--20\% depending on trauma type). The condition describes the development of anxiety, dissociative, and other symptoms within 1 month of a traumatic event, with required symptom clusters similar to PTSD (intrusion, negative mood, dissociation, avoidance, and arousal) but with an emphasis on dissociative symptoms. Mental health conditions affect thoughts, emotions, behaviors, and overall well-being. These conditions are among the most common health concerns worldwide, affecting people across all age groups. Mental health disorders result from complex interactions of biological, psychological, and social factors. Effective treatments are available, and recovery is possible with appropriate support.
\section{Causes}
Neurotransmitter imbalances (serotonin, dopamine, norepinephrine) play key roles in many mental health conditions. Genetic factors contribute to vulnerability, with heritability estimates varying by condition. Early life experiences, trauma, and chronic stress can alter brain development and function. Environmental factors including social support, socioeconomic status, and life events influence mental health.   
\section{Risk Factors}
\begin{itemize}[noitemsep]
  \item Genetic predisposition and family history
  \item Advancing age
  \item Lifestyle factors (diet, exercise, smoking, alcohol)
  \item Environmental exposures
  \item Underlying medical conditions
  \item Medication use
\end{itemize}
\section{Signs and Symptoms}
\subsection{Common symptoms}
\begin{itemize}[noitemsep]
  \item Clinical features last for 3--30 days after trauma
  \item If symptoms persist beyond 1 month, the diagnosis transitions to PTSD.
  \item Pain or discomfort in the affected area
  \end{itemize}
\subsection{Emergency warning signs}
\begin{itemize}[noitemsep]
  \item Severe or worsening symptoms of acute stress disorder require immediate medical evaluation
\end{itemize}
\section{Progression and Stages}
The condition may progress through different stages of severity over time. Early stages often involve mild or subtle symptoms that may be overlooked. Moderate stages involve more noticeable symptoms affecting daily function. Advanced stages may involve significant impairment and complications.
\section{Complications}
\begin{itemize}[noitemsep]
  \item Disease progression leading to worsening symptoms
  \item Secondary infections or injuries
  \item Side effects of treatment
  \item Reduced quality of life
  \item Emotional and psychological impact
\end{itemize}
\section{Diagnosis}
Doctors diagnose this using clinical interview.

\begin{itemize}[noitemsep]
  \item Comprehensive medical history and physical examination
  \item Laboratory tests to assess organ function
  \item Imaging studies to visualize affected structures
  \item Specialized testing depending on the affected system
  \item Consultation with relevant specialists
\end{itemize}
\section{Prevention}
\begin{itemize}[noitemsep]
  \item Maintain a healthy lifestyle with balanced diet and exercise
  \item Avoid known risk factors
  \item Attend regular medical check-ups for early detection
  \item Follow recommended screening guidelines
\end{itemize}
\section{Treatment Options}
Treatment options include trauma-focused CBT initiated within the first month. Psychotherapy (cognitive behavioral therapy, interpersonal therapy, psychodynamic therapy) Pharmacotherapy (antidepressants, antipsychotics, mood stabilizers, anxiolytics) Combined treatment approaches for optimal outcomes Hospitalization for acute crisis management Support groups and peer support programs Lifestyle interventions (exercise, sleep hygiene, stress management) Mindfulness and meditation practices Family therapy and psychoeducation Vocational and social rehabilitation Long-term maintenance therapy for chronic conditions

\begin{itemize}[noitemsep]
  \item Medications to manage symptoms and address underlying mechanisms
  \item Lifestyle modifications and self-care measures
  \item Physical therapy and rehabilitation
  \item Surgical intervention when indicated
  \item Regular monitoring and follow-up care
\end{itemize}
\section{Living with It}
\begin{itemize}[noitemsep]
  \item Follow your treatment plan as prescribed
  \item Maintain regular follow-up with your healthcare provider
  \item Make necessary lifestyle adjustments
  \item Seek support from family, friends, or support groups
  \item Stay informed about your condition
\end{itemize}
\section{Outlook}
\begin{itemize}[noitemsep]
  \item The outlook varies from person to person
  \item Approximately 50--80\% of patients with ASD develop PTSD within 6 months.
\end{itemize}
